\documentclass[10pt,a4paper]{article}
\usepackage[T1]{fontenc}
\usepackage[utf8]{inputenc}
\usepackage{lmodern}
\usepackage[margin=16mm]{geometry}
\usepackage{tabularx,booktabs}
\usepackage{enumitem}
\usepackage{titlesec}
\usepackage{textcomp}
\usepackage[hidelinks]{hyperref}

\setlength{\parindent}{0pt}
\setlength{\parskip}{3pt}
\setlist{nosep,leftmargin=*}
\titleformat{\section}{\large\bfseries}{\thesection.}{0.5em}{}
\titlespacing*{\section}{0pt}{7pt}{2pt}
\renewcommand{\arraystretch}{1.1}

\begin{document}

{\LARGE\bfseries Experimental Plan: Emotion-Contingent Navigation}\par
\vspace{2pt}
HCI Project, Semester 7 \textperiodcentered{} 25 October 2026\\
Running application: \url{https://github.com/VoidPsyker/HCI-Project}

\section{Research questions}

Online shops show every visitor the same menu, whatever mood they are in. We want to know
whether the best menu layout depends on how the shopper feels. If a frustrated shopper finds
a watch faster with a mega menu and a confused shopper does better with a sidebar, a shop
could switch its menu to suit the shopper's state.

We ask two questions. First, which of six menu structures helps people find products fastest
and most reliably in each of four emotional states? Second, can a model trained on clicks
and timing predict which state a shopper is in, well enough to choose the menu for them?

Hypothesis H1 says menu structure and emotional state interact, so the ranking of the menus
changes in at least some states. We do not decide in advance which menu suits which state;
the data decide. If one menu wins everywhere, that still answers the question: keep one
stable menu. Hypothesis H5 says the model will beat a majority-class baseline, though with
real uncertainty.

\section{The test application}

We built a fictional watch store, Meridian Timepieces, with 72 watches. A made-up brand
stops participants relying on what they know about real shops. All six menus show the same
4 departments, 12 subcategories, labels and products; only the structure changes. Search is
off during trials, so participants have to use the menu.

\begin{tabularx}{\linewidth}{@{}lX@{}}
\toprule
Menu & How it works \\
\midrule
A. Horizontal bar & Top bar, then a dropdown, then a side flyout (3 clicks) \\
B. Vertical sidebar & A category tree on the left that expands in place (3 clicks) \\
C. Mega menu & One large panel shows every subcategory and product type (2 clicks) \\
D. Fractal tiles & A $2\times2$ grid that zooms in level by level, with a breadcrumb (3 clicks) \\
E. Icon grid & Colour-coded rounded cards with icons at each level (3 clicks) \\
F. Radial lattice & Four compass sectors with deeper levels on outer rings (3 clicks) \\
\bottomrule
\end{tabularx}

\section{Emotional states and how we induce them}

Each block starts with a mild, 90-second activity.

\begin{tabularx}{\linewidth}{@{}lX@{}}
\toprule
State & Induction \\
\midrule
Happy & A short comedy clip plus the scenario ``you just found a great deal'' \\
Frustrated & A search box that returns ``No results'' for a valid query (three attempts at most) \\
Confused & Placing a watch into one of four overlapping categories, none of which fits \\
Bored & Checking the same \pounds240 price on 30 identical straps, with an option to stop early \\
\bottomrule
\end{tabularx}

Afterwards the participant rates how pleasant (1--9) and how alert (1--9) they feel and picks
the state that fits best. These self-reports tell us whether the induction worked and give
the model its labels. We never read emotion from faces or the webcam.

\section{Design and procedure}

Every participant uses all six menus in all four states (within-subjects). We counterbalance
state order with a $4\times4$ Latin square and menu order with a $6\times6$ Williams square,
and we rotate target products across six sets. The server assigns each sequence, so the
browser never picks a condition. Each watch has one finding task, such as ``Titanium dive
watch, 300\,m, deep blue dial'', rated easy, medium or hard by menu depth, how closely the
wording matches the label, and how many similar watches compete with it.

A pilot session takes about 90 minutes:
\begin{enumerate}
  \item Consent, separate eye-tracking consent, and screen and distance measurements.
  \item Demographics, a baseline mood check-in, instructions and six practice trials.
  \item Nine-point eye-tracking calibration. Error above 1.5\textdegree{} ends the session
        before scoring.
  \item Four blocks, each with an induction, a check-in, three trials, a second check-in with
        a fatigue rating, and a 60-second neutral reset.
  \item Questionnaires (SUS for each menu, NASA-TLX), a short interview and a debrief.
\end{enumerate}

A trial ends at the correct product page, at ``I can't find it'', or at 120 seconds. Feedback
stays neutral (``Response recorded'') so praise or error messages do not shift the mood.

\section{Data collection}

The app logs every action with a timestamp and the participant's pseudonymous ID, block,
trial, task, menu and assigned state. That covers each menu item opened or clicked (with
its screen position and level), back steps, wrong product pages, give-ups, check-ins and
questionnaire answers. The server stores everything in a database and also writes each
session's raw events to a separate backup file. Gaze samples go to a restricted store of
their own, and webcam video is never saved. The researcher exports CSV and JSONL files with
a checksum for each, so we can show later that no file changed. Rehearsal sessions with
synthetic data are flagged and excluded from analysis.

\begin{tabularx}{\linewidth}{@{}lX@{}}
\toprule
Type & What we measure \\
\midrule
Performance (per trial) & Success within 120\,s, completion time, wrong clicks, backtracking, path length \\
Emotional state & Pleasantness, alertness and dominant state after each induction and block \\
Experience & SUS for each menu, NASA-TLX at the end, open comments \\
Eye tracking & Fixations and dwell time on each menu area \\
\bottomrule
\end{tabularx}

\section{Model training: predicting the emotional state}

Study 1 gives us labelled behaviour. For each trial we compute features from the logs:
timing (time to first click, pauses, total time), clicking (click count, click rate, repeat
clicks on one item), navigation (back steps, category switches, menu openings, wrong
products, depth reached) and context (difficulty, menu, trial number). The context features
stop the model from learning the task instead of the mood. The label is the state the
participant reported (happy, frustrated, confused, bored or mixed), not the one we tried to
induce.

Our main model is an L2-penalised multinomial logistic regression with balanced class
weights. We chose it because its weights are readable and its probabilities can be
calibrated. Gradient boosting (LightGBM) serves as a check for non-linear gains. We validate
with leave-one-participant-out cross-validation, so the model is always tested on people it
never saw. We report macro-F1, balanced accuracy and the confusion matrix against a
majority-class baseline, and check calibration with the Brier score. H5 holds if macro-F1
beats the baseline by at least 0.10. We then freeze the model. In Study 3 it may switch the
menu only between tasks, only when it is at least 70\% confident, and the participant can
keep the current menu. A simulation over the Study 1 results will show how accurate the
model must be before switching beats a fixed menu. The model predicts what people report,
which is not the same as what they feel, and we will say so.

\section{Participants and phases}

The pilot has 8--10 participants with 12 scored trials each (3 per block). It tests the
method, not the hypotheses, and its data stay out of the main analysis. To move on, every
menu must reach 80\% practice success, at least 60\% of inductions must produce the intended
state with no distress rating above 6/10, overall success must fall between 75\% and 90\%,
and sessions must stay under 100 minutes with under 5\% missing data. Study 1 then runs 48
participants (72 if recruitment allows), aged 18--65 with normal or corrected vision, on 24
trials each: one per menu in every state. Studies 2 and 3, which test switching with known
and then predicted states, come after this semester.

\section{Analysis}

We model log completion time with a linear mixed-effects model and success with a logistic
one. Menu, state and their interaction are fixed effects; trial position and difficulty are
covariates; participant and target are random effects. A likelihood-ratio test on the
interaction tests H1, and we report effect sizes with 95\% confidence intervals. To pick a
menu for each state we fixed a rule in advance: keep menus within 5 percentage points of the
best success rate, then take the fastest. If no menu is clearly better, use the best overall.

\section{Ethics and risks}

We will not recruit until the ethics committee approves the study; the consent text and clips
in the app are placeholders until then. Participants can stop at any time, and a distress
rating above 6/10 sends them straight to the debrief. We never store names or webcam video,
and withdrawal deletes all of a participant's data. Mood carry-over is the main risk, which
the neutral reset and counterbalancing address. If sessions run long we will split Study 1
into two visits. Webcam eye tracking ran at about 48 frames per second on our laptop, just
under the 50\,fps target, so we will retest it on the lab machine.

\end{document}
